\documentclass[10pt,a4paper]{article}

\usepackage[T1]{fontenc}
\usepackage[utf8]{inputenc}
\usepackage{lmodern}
\usepackage{microtype}
\usepackage{geometry}
\geometry{left=18mm,right=18mm,top=18mm,bottom=20mm,headheight=20pt}
\usepackage{setspace}
\setstretch{1.02}
\usepackage{amsmath,amssymb,mathtools,bm}
\usepackage{array,tabularx,booktabs}
\usepackage{enumitem}
\usepackage{xcolor}
\usepackage[most]{tcolorbox}
\usepackage{needspace}
\usepackage{etoolbox}
\usepackage{fancyhdr}
\usepackage{titlesec}
\usepackage[hidelinks]{hyperref}

\definecolor{ink}{HTML}{414141}
\definecolor{vuBurgundy}{HTML}{78003F}
\definecolor{vuPink}{HTML}{E64164}
\definecolor{blue}{HTML}{174F78}
\definecolor{paleBlue}{HTML}{F1F6F9}
\definecolor{paleBurgundy}{HTML}{FAF2F6}
\definecolor{palePink}{HTML}{FDF3F5}
\definecolor{softGray}{HTML}{F5F5F5}
\definecolor{workGray}{HTML}{ABB3BB}

\color{ink}
\setlength{\parindent}{0pt}
\setlength{\parskip}{0.42em}
\raggedbottom
\clubpenalty=10000
\widowpenalty=10000
\displaywidowpenalty=10000
\setlist[itemize]{leftmargin=1.5em,itemsep=0.18em,topsep=0.20em}
\setlist[enumerate]{leftmargin=3.0em,itemsep=0.20em,topsep=0.20em}

\titleformat{\section}
  {\Large\bfseries\color{vuBurgundy}}{\thesection}{0.65em}{}
\titleformat{\subsection}
  {\large\bfseries\color{blue}}{\thesubsection}{0.65em}{}
\titlespacing*{\section}{0pt}{0.9em}{0.50em}
\titlespacing*{\subsection}{0pt}{0.7em}{0.35em}
\pretocmd{\section}{\Needspace{14\baselineskip}}{}{}
\pretocmd{\subsection}{\Needspace{5\baselineskip}}{}{}

\pagestyle{fancy}
\fancyhf{}
\fancyhead[L]{\small\color{vuBurgundy}\textbf{Maps and isomorphisms}}
\fancyhead[R]{\small\color{ink!65}Functions and structure}
\fancyfoot[C]{\small\color{vuBurgundy}\thepage}
\renewcommand{\headrulewidth}{0.4pt}
\renewcommand{\headrule}{\hbox to\headwidth{%
  \color{vuBurgundy}\leaders\hrule height \headrulewidth\hfill}%
  \vskip-\headrulewidth}

\hypersetup{
  pdftitle={Maps and Isomorphisms},
  pdfsubject={Injective and surjective maps, isomorphisms, diffeomorphisms, and proof strategies},
  colorlinks=true,
  linkcolor=vuBurgundy,
  urlcolor=blue
}

\newcommand{\R}{\mathbb{R}}
\newcommand{\C}{\mathbb{C}}
\newcommand{\Z}{\mathbb{Z}}
\newcommand{\Q}{\mathbb{Q}}
\newcommand{\N}{\mathbb{N}}
\newcommand{\End}{\operatorname{End}}

\newtcbtheorem[number within=section]{definition}{Definition}{
  enhanced,colback=paleBlue,colframe=blue,
  colbacktitle=blue,coltitle=white,fonttitle=\bfseries,
  boxrule=0.7pt,arc=1.4mm,
  left=2.5mm,right=2.5mm,top=1.5mm,bottom=1.5mm,
  before skip=5pt,after skip=6pt
}{def}

\newtcbtheorem[use counter from=definition]{proposition}{Proposition}{
  enhanced,colback=paleBurgundy,colframe=vuBurgundy,
  colbacktitle=vuBurgundy,coltitle=white,fonttitle=\bfseries,
  boxrule=0.7pt,arc=1.4mm,
  left=2.5mm,right=2.5mm,top=1.5mm,bottom=1.5mm,
  before skip=5pt,after skip=6pt
}{prop}

\newtcbtheorem[use counter from=definition]{examplebox}{Example}{
  enhanced,colback=softGray,colframe=ink!45,
  colbacktitle=ink!65,coltitle=white,fonttitle=\bfseries,
  boxrule=0.6pt,arc=1.4mm,
  left=2.5mm,right=2.5mm,top=1.5mm,bottom=1.5mm,
  before skip=5pt,after skip=6pt
}{ex}

\newtcolorbox{precision}[1][Remark]{
  enhanced,colback=palePink,colframe=vuPink,
  title={#1},colbacktitle=vuBurgundy,coltitle=white,fonttitle=\bfseries,
  boxrule=0.7pt,arc=1.4mm,
  left=2.5mm,right=2.5mm,top=1.5mm,bottom=1.5mm,
  before skip=5pt,after skip=6pt
}

\newtcolorbox{summarybox}[1]{
  enhanced,colback=softGray,colframe=workGray,
  title={#1},colbacktitle=ink!70,coltitle=white,fonttitle=\bfseries,
  boxrule=0.65pt,arc=1.4mm,
  left=2.5mm,right=2.5mm,top=1.5mm,bottom=1.5mm,
  before skip=5pt,after skip=6pt
}

\newcommand{\keyidea}[1]{%
  \par\smallskip
  \begin{minipage}{\linewidth}
  \noindent\textbf{\color{vuBurgundy}Key idea.} #1
  \end{minipage}
  \par\smallskip
}

\BeforeBeginEnvironment{proposition}{\Needspace{12\baselineskip}}

\begin{document}

\begin{center}
  {\Huge\bfseries\color{vuBurgundy}Maps and Isomorphisms\par}
  \vspace{2mm}
  {{\large Functions, algebraic structure, and smooth maps\par}}
\end{center}

The properties of a map depend on its domain, its codomain, and the structure
that it preserves. Injectivity and surjectivity determine whether a map has an
inverse as a function. An isomorphism must also preserve the relevant
operations. We develop these notions and their proof methods, then apply them
to the identification of vectors in $\R^n$ with derivations at a point.

\section{Maps, images, and preimages}

\begin{definition}{Function, domain, and codomain}{function}
A function $F:A\to B$ assigns to each $a\in A$ exactly one element
$F(a)\in B$. The set $A$ is the \emph{domain}, and $B$ is the \emph{codomain}.
The \emph{image} is the subset
\[
  F(A)=\{F(a):a\in A\}\subseteq B.
\]
For $b\in B$, the set $F^{-1}(\{b\})=\{a\in A:F(a)=b\}$ consists of the
preimages of $b$. This notation for a preimage set does not require an inverse
function.
\end{definition}

A formula alone does not determine the properties of a map. For example,
$x\mapsto x^2$ behaves differently as a map $\R\to\R$, as a map
$[0,\infty)\to\R$, and as a map $[0,\infty)\to[0,\infty)$. The domain and codomain must therefore be specified as part of the map.

\begin{summarybox}{Logical forms}
\renewcommand{\arraystretch}{1.35}
\begin{tabularx}{\linewidth}{>{\bfseries}p{2.5cm}X}
Injective & For all $a,a'\in A$, $F(a)=F(a')$ implies $a=a'$.\\
Surjective & For every $b\in B$, there exists $a\in A$ with $F(a)=b$.\\
Bijective & Both statements hold.\\
\end{tabularx}
\end{summarybox}

The quantifiers determine the initial step of a proof. Injectivity concerns
two arbitrary elements of the domain. Surjectivity concerns an arbitrary
element of the codomain and the existence of a preimage.

\begin{precision}[Maps on equivalence classes]
If an input is an equivalence class, a formula involving a representative
might give different outputs for different representatives. Before studying
injectivity or surjectivity, prove that the output is independent of the
representative. For example, $[f]_p\mapsto f(p)$ is well defined on smooth
function germs because functions that agree near $p$ have the same value at
$p$. This is a separate question from whether the map is injective.
\end{precision}

\section{Injective maps}

\begin{definition}{Injective map}{injective}
A map $F:A\to B$ is \emph{injective} if
\[
  F(a)=F(a')\quad\Longrightarrow\quad a=a'
  \qquad\text{for every }a,a'\in A.
\]
Equivalently, no element of $B$ has two distinct preimages in $A$.
\end{definition}

\begin{summarybox}{Establishing injectivity}
Take arbitrary $a,a'\in A$ and assume $F(a)=F(a')$. Use that equality and
the definition of $F$ to derive $a=a'$. You may instead assume $F(a)=F(a')$
and show that $a-a'=0$ when $A$ is a vector space. To disprove injectivity,
find distinct $a,a'$ with the same image.
\end{summarybox}

\begin{examplebox}{Restriction of the square map}{squares-inj}
The map $q:\R\to\R$, $q(x)=x^2$, is not injective because
$q(1)=q(-1)=1$ but $1\ne-1$.

The restriction $q:[0,\infty)\to\R$ is injective. If $x,y\geq0$ and
$x^2=y^2$, then $x=y$ because nonnegative numbers have a unique
nonnegative square root. The conclusion depends on the domain.
\end{examplebox}

\begin{proposition}{The kernel criterion}{kernel}
Let $T:V\to W$ be linear. Then $T$ is injective if and only if
$\ker T=\{0_V\}$, where
$\ker T=\{v\in V:T(v)=0_W\}$.
\end{proposition}

\textit{Proof.} If $T$ is injective and $T(v)=0_W=T(0_V)$, then $v=0_V$.
Conversely, if $\ker T=\{0_V\}$ and $T(u)=T(v)$, linearity gives
$T(u-v)=T(u)-T(v)=0_W$. Hence $u-v=0_V$, so $u=v$.
\hfill$\square$

For a linear map, injectivity can therefore be established by fixing
$v\in V$ and assuming $T(v)=0$. If $T(v)$ is an operator, this equality
means that it vanishes on every function in its domain. Suitable test
functions may be chosen to recover the components of $v$.

\section{Surjective maps}

\begin{definition}{Surjective map}{surjective}
A map $F:A\to B$ is \emph{surjective} if $F(A)=B$. Written with quantifiers,
this says
\[
  \text{for every }b\in B\text{ there exists }a\in A\text{ such that }F(a)=b.
\]
Different inputs may have the same output. Surjectivity requires at least one
preimage for each target, not exactly one.
\end{definition}

\begin{summarybox}{Establishing surjectivity}
Take an arbitrary $b\in B$. Solve $F(a)=b$ or construct a suitable $a$.
Check that $a\in A$, then compute $F(a)=b$. To disprove surjectivity, exhibit
one $b\in B$ that has no preimage.
\end{summarybox}

\begin{examplebox}{The codomain of the square map}{squares-surj}
The map $q:\R\to\R$, $q(x)=x^2$, is not surjective because $-1$ is not
a square of a real number. The map $q:\R\to[0,\infty)$ with the same rule
is surjective: given any $b\geq0$, take $a=\sqrt b\in\R$ and obtain
$q(a)=b$. The codomain determines whether the map is surjective.
\end{examplebox}

\begin{examplebox}{Evaluation on smooth functions}{evaluation}
Fix $p\in U$, where $U\subseteq\R^n$ is open. Consider
\[
  \operatorname{ev}_p:C^\infty(U)\to\R,\qquad f\mapsto f(p).
\]
It is surjective. Given $c\in\R$, choose the constant smooth function
$f(x)=c$. Then $\operatorname{ev}_p(f)=c$. It is not injective: the zero
function and $x\mapsto x^1-p^1$ are distinct but both take value zero at
$p$. This example also shows why a surjective homomorphism need not be an
isomorphism.
\end{examplebox}

\begin{precision}[An arbitrary element of the codomain]
Writing ``let $b=F(a)$'' merely selects an element already in the image.
Surjectivity requires the opposite order: fix an arbitrary $b\in B$ and
produce an $a\in A$ for which $F(a)=b$.
\end{precision}

\section{Bijective maps and inverses}

\begin{definition}{Bijection and inverse}{bijection}
A map $F:A\to B$ is \emph{bijective} if it is injective and surjective.
Then every $b\in B$ has exactly one preimage $a\in A$. The \emph{inverse
map} $F^{-1}:B\to A$ assigns that unique $a$ to $b$.
\end{definition}

\begin{proposition}{Criterion for an inverse map}{inverse-criterion}
A map $F:A\to B$ is bijective if and only if there is a map $G:B\to A$
for which
\[
  G\circ F=\operatorname{id}_A,
  \qquad F\circ G=\operatorname{id}_B.
\]
In that case $G=F^{-1}$.
\end{proposition}

\textit{Proof.} If $F$ is bijective, its inverse satisfies both equations.
Conversely, if $F(a)=F(a')$, apply $G$ to obtain $a=a'$, so $F$ is
injective. Given $b\in B$, let $a=G(b)$. Then $F(a)=F(G(b))=b$, so $F$ is
surjective.\hfill$\square$

One equation by itself has a more limited consequence. A \emph{left inverse}
$G\circ F=\operatorname{id}_A$ proves that $F$ is injective. A
\emph{right inverse} $F\circ G=\operatorname{id}_B$ proves that $F$ is
surjective. These are distinct conditions on the compositions.

\begin{examplebox}{The cubic map}{cubic}
Let $F:\R\to\R$ be $F(x)=x^3$. Given $y\in\R$, solving $y=x^3$ gives
$x=\sqrt[3]{y}$. The map $G(y)=\sqrt[3]{y}$ satisfies
$G(F(x))=x$ and $F(G(y))=y$. Hence $F$ is bijective and $G=F^{-1}$.
\end{examplebox}

\begin{summarybox}{The square map with different domains and codomains}
\renewcommand{\arraystretch}{1.25}
\begin{tabularx}{\linewidth}{>{\raggedright\arraybackslash}p{4.8cm}c c X}
\toprule
Map $x\mapsto x^2$ & Injective? & Surjective? & Reason\\
\midrule
$\R\to\R$ & no & no & $1$ and $-1$ have the same image, and negatives are missed\\
$\R\to[0,\infty)$ & no & yes & each target has a square root\\
$[0,\infty)\to\R$ & yes & no & unique nonnegative root, but negatives are missed\\
$[0,\infty)\to[0,\infty)$ & yes & yes & inverse $y\mapsto\sqrt y$\\
\bottomrule
\end{tabularx}
\end{summarybox}

\section{Homomorphisms and isomorphisms}

Injectivity and surjectivity discuss the correspondence between inputs and
outputs. A \emph{homomorphism} must also respect operations. Which operations
it must preserve depends on the structures involved.

\begin{definition}{Homomorphisms}{homomorphism}
For groups $(G,\cdot)$ and $(H,\star)$, a group homomorphism $h:G\to H$
satisfies $h(ab)=h(a)\star h(b)$. For vector spaces over $K$, a
homomorphism is a $K$-linear map:
\[
  T(u+v)=T(u)+T(v),\qquad T(ru)=rT(u).
\]
For unital $K$-algebras, a unital algebra homomorphism is $K$-linear and
also satisfies $h(ab)=h(a)h(b)$ and $h(1_A)=1_B$.
\end{definition}

The homomorphism condition is checked separately for every operation
required by the given structures. A map may preserve one operation and fail
to preserve another. For example, a set bijection between vector spaces need not be
linear: $F:\R\to\R$, $F(x)=x^3$, is bijective, but
$F(1+1)=8\ne2=F(1)+F(1)$.

\begin{definition}{Isomorphism}{isomorphism}
An \emph{isomorphism} of a specified kind of algebraic structure is a
bijective homomorphism of that structure. A linear isomorphism is a bijective
linear map. A unital algebra isomorphism is a bijective unital algebra
homomorphism. The inverse of an algebraic isomorphism preserves the same
operations.
\end{definition}

\begin{proposition}{Inverse of a linear isomorphism}{inverse-linear}
If $T:V\to W$ is a bijective linear map, then $T^{-1}:W\to V$ is linear.
\end{proposition}

\textit{Proof.} Take arbitrary $w_1,w_2\in W$ and $r,s\in K$. Since $T$ is
surjective, write $w_1=T(v_1)$ and $w_2=T(v_2)$. Then
\[
  T^{-1}(rw_1+sw_2)
  =T^{-1}\bigl(T(rv_1+sv_2)\bigr)
  =rv_1+sv_2
  =rT^{-1}(w_1)+sT^{-1}(w_2).
\]
The same idea proves preservation of products for a bijective algebra
homomorphism.\hfill$\square$

\begin{summarybox}{Conditions for the principal types of maps}
\renewcommand{\arraystretch}{1.28}
\begin{tabularx}{\linewidth}{>{\bfseries}p{4.6cm}X}
Bijection of sets & Injectivity and surjectivity.\\
Linear isomorphism & Linearity, injectivity, and surjectivity.\\
Algebra isomorphism & Required algebra laws, injectivity, and surjectivity.\\
Diffeomorphism & Smoothness, bijectivity, and smoothness of the inverse.\\
\end{tabularx}
\end{summarybox}

The word \emph{isomorphic} always needs an understood structure. Two sets can
be bijective without being isomorphic as groups, vector spaces, or algebras.
An \emph{automorphism} is an isomorphism from a structure to itself.

\section{Homeomorphisms and diffeomorphisms}

\begin{definition}{Homeomorphism and diffeomorphism}{diffeomorphism}
A \emph{homeomorphism} $F:X\to Y$ is a bijective continuous map whose
inverse is continuous. A \emph{diffeomorphism} $F:M\to N$ between smooth
manifolds is a bijective smooth map whose inverse is smooth. For open subsets
of Euclidean space, smoothness means that derivatives of all orders exist and
are continuous in local coordinates.
\end{definition}

These terms describe which structure is preserved. A diffeomorphism is also a
homeomorphism, since smooth maps are continuous. Bijectivity by itself says
nothing about continuity or differentiability.

\begin{examplebox}{A smooth bijection without a smooth inverse}{cubic-diff}
The map $F:\R\to\R$, $F(x)=x^3$, is smooth and bijective. Its inverse
$F^{-1}(y)=\sqrt[3]{y}$ is continuous, so $F$ is a homeomorphism. But
\[
  \frac{F^{-1}(h)-F^{-1}(0)}{h}
  =\frac{\sqrt[3]{h}}{h}=|h|^{-2/3}\quad(h\ne0),
\]
which has no finite limit at $0$. Thus $F^{-1}$ is not differentiable at
$0$, and $F$ is not a diffeomorphism.
\end{examplebox}

\begin{examplebox}{The exponential diffeomorphism}{exp-diff}
The map $E:\R\to(0,\infty)$, $E(x)=e^x$, is smooth. The function
$L:(0,\infty)\to\R$, $L(y)=\log y$, is smooth, and
$L(E(x))=x$, $E(L(y))=y$. Thus $E$ is a diffeomorphism. If its codomain
were $\R$, it would fail to be surjective.
\end{examplebox}

\begin{precision}[Local and global inverses]
For a smooth map between open subsets of $\R^n$, an invertible derivative
at one point gives a smooth inverse \emph{near that point} by the inverse
function theorem. It does not establish a global inverse. To prove a global
diffeomorphism, also establish global bijectivity and smoothness of the
inverse, either explicitly or by combining bijectivity with local inverses.
\end{precision}

\section{Vectors as point derivations}

We now prove the identification between vectors and point derivations. Fix
$p\in U\subseteq\R^n$, with $U$ open. Write $C_p^\infty$ for the
algebra of germs at $p$ of smooth real-valued functions near $p$.

\begin{definition}{Point derivation}{point-derivation}
A \emph{derivation at $p$} is an $\R$-linear map
$D:C_p^\infty\to\R$ satisfying
\[
  D([fg]_p)=D([f]_p)g(p)+f(p)D([g]_p).
\]
Let $\operatorname{Der}_p$ be the real vector space of all such derivations,
with pointwise addition and scalar multiplication of maps.
\end{definition}

For $v=(v^1,\ldots,v^n)\in\R^n$, define
\[
  \Phi:\R^n\longrightarrow\operatorname{Der}_p,
  \qquad \Phi(v)=D_v,
  \qquad D_v([f]_p)=\sum_{i=1}^n v^i\,\partial_i f(p).
\]
We must prove that $\Phi$ is well defined and linear, then injective, then
surjective. The vector $v$ is fixed when $D_v$ acts on a function germ
$[f]_p$.

\begin{proposition}{Linearity of $\Phi$}{phi-well}
The displayed formula defines a point derivation for every $v\in\R^n$,
and $\Phi$ is $\R$-linear.
\end{proposition}

\textit{Proof.} If two representatives have the same germ at $p$, they agree
on a neighbourhood of $p$. Their partial derivatives at $p$ therefore agree,
so $D_v([f]_p)$ is independent of the representative. Linearity in $[f]_p$
follows from linearity of partial derivatives. The product rule gives
\[
  D_v([fg]_p)
  =\sum_i v^i\bigl(\partial_i f(p)g(p)+f(p)\partial_i g(p)\bigr)
  =D_v([f]_p)g(p)+f(p)D_v([g]_p).
\]
Finally, for $v,w\in\R^n$ and $a,b\in\R$, evaluating on any $[f]_p$
shows $D_{av+bw}([f]_p)=aD_v([f]_p)+bD_w([f]_p)$.
Hence $\Phi(av+bw)=a\Phi(v)+b\Phi(w)$.\hfill$\square$

\begin{proposition}{Injectivity of $\Phi$}{phi-inj}
The map $\Phi$ is injective.
\end{proposition}

\textit{Proof.} Since $\Phi$ is linear, it suffices to show that its kernel
is zero. Fix $v\in\R^n$ and suppose $\Phi(v)=D_v=0$, the zero derivation.
The equality $D_v=0$ means that $D_v([f]_p)=0$ for \emph{every} germ, so we
may choose a useful one. For each $j$, let $x^j$ denote the $j$th coordinate
function. Then $\partial_i x^j(p)=\delta_i^{\,j}$, and
\[
  0=D_v([x^j]_p)
   =\sum_{i=1}^n v^i\delta_i^{\,j}=v^j.
\]
Every component of $v$ is zero. Thus $v=0$, so $\ker\Phi=\{0\}$.
\hfill$\square$

The coordinate functions need not represent arbitrary germs. The hypothesis
$D_v=0$ asserts that the operator vanishes on every germ. Evaluation on the
coordinate germs isolates the components of the fixed vector $v$.

\begin{proposition}{Surjectivity of $\Phi$}{phi-surj}
Every derivation $D\in\operatorname{Der}_p$ equals $D_v$ for a unique
$v\in\R^n$. Hence $\Phi$ is surjective.
\end{proposition}

\textit{Proof.} Let $D\in\operatorname{Der}_p$ be arbitrary. Define a vector $v$ by
\[
  v^i=D([x^i]_p),\qquad i=1,\ldots,n.
\]
We must show $D([f]_p)=D_v([f]_p)$ for \emph{every} smooth germ.
The Leibniz rule gives $D([1]_p)=D([1]_p[1]_p)=2D([1]_p)$, so
$D([1]_p)=0$. By $\R$-linearity, $D$ kills every constant germ.

Choose a small convex ball around $p$ on which $f$ is defined. On that ball,
\[
  f(x)-f(p)=\sum_{i=1}^n (x^i-p^i)g_i(x),
  \qquad
  g_i(x)=\int_0^1\partial_i f\bigl(p+t(x-p)\bigr)\,dt.
\]
This identity follows by applying the fundamental theorem of calculus to
$t\mapsto f(p+t(x-p))$. Each $g_i$ is smooth near $p$, and
$g_i(p)=\partial_i f(p)$. Passing to germs and applying $D$ gives
\begin{align*}
  D([f]_p)
   &=\sum_i D\bigl([x^i-p^i]_p[g_i]_p\bigr)\\
   &=\sum_i\left(D([x^i-p^i]_p)g_i(p)
       +(x^i(p)-p^i)D([g_i]_p)\right)\\
   &=\sum_i D([x^i]_p)\partial_i f(p)
    =\sum_i v^i\partial_i f(p)=D_v([f]_p).
\end{align*}
Thus $D=D_v=\Phi(v)$ as maps on all germs. Uniqueness of $v$ follows from
injectivity.\hfill$\square$

The injectivity argument recovers a vector from the values of its
derivation on coordinate functions. The surjectivity argument constructs a
vector from an arbitrary derivation and proves equality of the resulting
operators on every germ.

The three propositions prove that $\Phi:\R^n\to\operatorname{Der}_p$ is a
linear isomorphism. Its inverse is concrete:
\[
  \Phi^{-1}(D)=\bigl(D([x^1]_p),\ldots,D([x^n]_p)\bigr).
\]

\section{Proof methods and exercises}

\subsection{Summary of the proof criteria}

\begin{summarybox}{Standard proof steps}
\begin{enumerate}[label=\arabic*.]
  \item State the domain and codomain. If inputs are equivalence classes,
  check well-definedness first.
  \item For injectivity, take two inputs with equal images, or take an
  element of the kernel when the map is linear.
  \item For surjectivity, take an arbitrary target and construct a preimage.
  \item For a homomorphism, verify the required operations. For a
  diffeomorphism, verify smoothness in both directions.
\end{enumerate}
\end{summarybox}

\begin{precision}[Insufficient arguments]
\begin{itemize}
  \item Checking $F(a)\ne F(a')$ for one pair does not prove injectivity.
  The claim is about \emph{all} distinct pairs.
  \item Producing one output $F(a)$ does not prove surjectivity. The target
  must be arbitrary in the stated codomain.
  \item A bijective map of sets is not automatically an isomorphism of
  vector spaces or a diffeomorphism. Check the relevant structure.
\end{itemize}
\end{precision}

For smooth vector fields, the following pointwise operations will be used. If $X$ is a smooth
vector field and $f$ is smooth, then
$(Xf)(p)=X_p(f)=\sum_i X^i(p)\partial_i f(p)$ defines a smooth
\emph{function}. For a smooth function $h$, the field $hX$ has value
$(hX)_p=h(p)X_p$.

The following exercises apply the preceding criteria. In each case, the
domain and codomain form part of the statement to be proved.

\subsection{Exercises}

\begin{enumerate}[label=\textbf{\arabic*.},leftmargin=2.1em,itemsep=0.5em]
  \item Let $T:\R^2\to\R^2$ be $T(x,y)=(x+y,x-y)$. Prove that $T$ is a
  linear isomorphism and find its inverse.
  \item Let $P:\R^3\to\R^2$ be $P(x,y,z)=(x+z,y)$. Determine whether it is
  injective and surjective. Give witnesses.
  \item Consider $\operatorname{ev}_p:C^\infty(U)\to\R$. Is it an
  isomorphism of real vector spaces? Explain.
  \item Let $F:\R\to\R$, $F(x)=x^3$. Which of ``bijection'',
  ``homeomorphism'', and ``diffeomorphism'' apply?
  \item Suppose $F:A\to B$ has a left inverse $G:B\to A$. Prove
  injectivity directly from $F(a)=F(a')$. Does this alone prove
  surjectivity? Give an example.
  \item Fix $p\in U$. Show that $J:C_p^\infty\to\R$ given by
  $J([f]_p)=\partial_1 f(p)$ is well defined, linear, and surjective.
  Is it injective?
  \item For $n\geq2$, consider $Q:\Z\to\Z/n\Z$, $k\mapsto[k]$.
  Prove that it is an additive group homomorphism and surjective. Is it
  injective? Is it an isomorphism of groups?
  \item For a smooth vector field $X$ and smooth functions $f,h$ on $U$,
  prove $(hX)f=h(Xf)$ by evaluating at an arbitrary point. Which of
  $hX$ and $Xf$ is a vector field?
\end{enumerate}

\subsection{Solutions}

\textbf{1.} Componentwise calculation gives
$T(a(x,y)+b(u,v))=aT(x,y)+bT(u,v)$, so $T$ is linear. Given
$(s,t)\in\R^2$, solve $s=x+y$, $t=x-y$. The unique solution is
$x=(s+t)/2$, $y=(s-t)/2$. Thus
$T^{-1}(s,t)=((s+t)/2,(s-t)/2)$, and the formulas compose to the identity
in both orders. Therefore $T$ is a linear isomorphism.

\textbf{2.} The nonzero vector $(1,0,-1)$ lies in $\ker P$, so $P$ is not
injective. Given $(a,b)\in\R^2$, take $(x,y,z)=(a,b,0)$. Its image is
$(a,b)$. Hence $P$ is surjective.

\textbf{3.} Evaluation is linear and surjective by constant functions. It
is not injective: $f(x)=x^1-p^1$ is a nonzero smooth function with
$f(p)=0$. Therefore it is not a vector-space isomorphism.

\textbf{4.} The cubic is bijective with continuous inverse
$y\mapsto\sqrt[3]{y}$, so it is a homeomorphism. Its inverse is not
differentiable at $0$, so it is not a diffeomorphism.

\textbf{5.} If $F(a)=F(a')$, then
$a=G(F(a))=G(F(a'))=a'$. This proves injectivity. The inclusion
$F:\{0\}\to\{0,1\}$ has a left inverse $G:\{0,1\}\to\{0\}$, but $F$
misses $1$, so a left inverse alone does not prove surjectivity.

\textbf{6.} Equal germs agree on a neighbourhood of $p$, so their first
partial derivatives at $p$ agree. This proves well-definedness. The map is
linear by the derivative rules. Given $c\in\R$, choose $f(x)=cx^1$. Then $J([f]_p)=c$. The nonzero constant germ $[1]_p$ lies in the kernel,
so $J$ is not injective.

\textbf{7.} Since $[k+\ell]=[k]+[\ell]$, the map is a group homomorphism.
Given a residue class $[r]$, the integer $r$ maps to it, so $Q$ is
surjective. But $Q(0)=Q(n)=[0]$ and $0\ne n$. It is not injective and
therefore not a group isomorphism.

\textbf{8.} At any $p\in U$, the field $hX$ has value
$(hX)_p=h(p)X_p$. Hence
$((hX)f)(p)=(hX)_p(f)=h(p)X_p(f)=(h(Xf))(p)$.
Equality at every $p$ proves $(hX)f=h(Xf)$ as functions.
The object $hX$ is a vector field. The object $Xf$ is a smooth function.

\end{document}
